\section{Das konstruktible Universum}

\subsection{Definierbare Teilmengen und die Hierarchie}

\FormulaDefDeltaK[Konstruktible Hierarchie]
{\begin{aligned}
 \BXLVIIIDef(X)
 &=\{\{z\in X:(X,\in)\models\varphi(z,\bar a)\}:
               \varphi\text{ Formel},\ \bar a\in X^{<\omega}\},\\
 L_0&=\varnothing,\qquad
 L_{\alpha+1}=\BXLVIIIDef(L_\alpha),\\
 L_\lambda&=\bigcup_{\alpha<\lambda}L_\alpha,\qquad
 x\in L\ \Longleftrightarrow\ \exists\alpha\in\BXLVIIIOrd\ (x\in L_\alpha)
\end{aligned}}
{B48LDefinition}{
  \DeltaRow{Menge}{X}
  \DeltaRow{Limesordinal}{\lambda}
}
\noindent
Die Erfüllungsrelation wird hier nur für die Mengenstruktur
\((X,\in)\) verwendet. Es gibt eine Menge von Formelcodes und endlichen
Parametertupeln; Aussonderung und Ersetzung machen
\(\BXLVIIIDef(X)\) zu einer Menge. Die Hierarchie entsteht durch den
bereits bewiesenen Rekursionssatz. Eine Formel \(\varphi^L\) entsteht
aus \(\varphi\), indem jeder Quantor auf die definierbare Klasse \(L\)
eingeschränkt wird.

\FormulaThmDeltaK[Grundlegende Eigenschaften der \(L\)-Stufen]
{\begin{aligned}
 \alpha\le\beta&\ \Longrightarrow\ L_\alpha\subseteq L_\beta,\\
 L_\alpha&\text{ ist transitiv},\qquad
 L_\alpha\cap\BXLVIIIOrd=\alpha,\\
 \alpha\ge\omega&\ \Longrightarrow\ |L_\alpha|=|\alpha|
\end{aligned}}
{B48LStages}{
  \DeltaRow{Ordinale}{\alpha,\beta}
}
\begin{tabproof}
 \proofstep{}{\BXLVIIIText{Ist \(L_\alpha\) transitiv und
 \(a\in L_\alpha\), so ist
 \(a=\{z\in L_\alpha:z\in a\}\) definierbar.
 Deshalb \(L_\alpha\subseteq\BXLVIIIDef(L_\alpha)\).
 Jedes Element von \(\BXLVIIIDef(L_\alpha)\) ist eine Teilmenge von
 \(L_\alpha\), womit auch die nächste Stufe transitiv ist.}}
 {\BXLVIIIWhy{Induktion: Nachfolgerschritt}}
 \proofstep{}{\BXLVIIIText{An Limesstellen erhalten Vereinigungen
 aufsteigender transitiver Mengen diese Eigenschaften.
 Beginnend mit der leeren Menge folgen Transitivität und Monotonie
 für alle Stufen.}}
 {\BXLVIIIWhy{Transfinite Induktion}}
 \proofstep{}{\BXLVIIIText{Die Ordinale in einer transitiven Stufe
 sind intern durch die beschränkten Bedingungen
 \enquote{transitiv und durch \(\in\) strikt linear geordnet}
 erkennbar; Fundierung macht daraus eine Wohlordnung.
 Unter \(L_\alpha\cap\BXLVIIIOrd=\alpha\) ist daher \(\alpha\)
 eine definierbare Teilmenge von \(L_\alpha\), also Element der
 nächsten Stufe. Keine größere Ordinalzahl kann Teilmenge von
 \(L_\alpha\) sein. Der Limesfall ist die Vereinigung.}}
 {\BXLVIIIWhy{Induktion über die Ordinalteile}}
 \proofstep{}{\BXLVIIIText{Für unendliches \(X\) gibt es höchstens
 \(\aleph_0\cdot|X|^{<\omega}=|X|\) definierbare Teilmengen.
 Die endlichen Anfangsstufen sind endlich; \(L_\omega\) ist
 abzählbar. Bei Nachfolgern wächst die unendliche Mächtigkeit
 nicht, bei Limiten ist
 \(|L_\lambda|\le|\lambda|\cdot|\lambda|=|\lambda|\).}}
 {\BXLVIIIWhy{Codierung durch Formeln und Parameter; Kardinalprodukte}}
 \proofstep{}{|L_\alpha|=|\alpha|\quad(\alpha\ge\omega)}
 {\BXLVIIIWhy{Obere Schranke aus 4, untere aus \(\alpha\subseteq L_\alpha\)}}
\end{tabproof}

\subsection{Reflexion und die ZF-Axiome}

\FormulaThmDeltaK[Reflexion einer endlichen Formelliste in \(L\)]
{\begin{gathered}
 \mathcal F\text{ endliche Formelliste},\quad\gamma\in\BXLVIIIOrd\\
 \vdash\ \exists\theta>\gamma\quad
 \forall\varphi\in\mathcal F\ \forall\bar a\in L_\theta\
 \bigl(\varphi^L(\bar a)\longleftrightarrow
                L_\theta\models\varphi(\bar a)\bigr)
 \end{gathered}}
{B48LReflection}{
  \DeltaRow{Parameter}{\gamma,\mathcal F}
}
\begin{tabproof}
 \proofstep{1}{\mathcal F\text{ endlich},\quad\gamma\in\BXLVIIIOrd}{\rA}
 \proofstep{1}{\BXLVIIIText{Schließe die Liste unter Teilformeln ab.
 Für jede Existenzteilformel und jedes Parametertupel aus einer
 festen Stufe \(L_\alpha\), für das ein Zeuge in \(L\) existiert,
 gibt es eine kleinste Stufe, welche einen Zeugen enthält.}}
 {\BXLVIIIWhy{Ordinalwohlordnung; feste relativierte Formeln}}
 \proofstep{1}{\BXLVIIIText{Ersetzung sammelt diese kleinsten
 Ordinale, da es nur eine Menge von Tupeln und endlich viele
 Formeln gibt. Wähle eine Ordinalzahl \(h(\alpha)>\alpha\),
 deren Stufe alle so erfassten Zeugen enthält.}}
 {\BXLVIIIWhy{Ersetzung und Supremum}}
 \proofstep{1}{\alpha_0>\gamma,\quad
 \alpha_{n+1}=h(\alpha_n),\quad
 \theta=\sup_{n<\omega}\alpha_n}
 {\BXLVIIIWhy{Abzählbare Rekursion}}
 \proofstep{1}{\BXLVIIIText{Jedes endliche Tupel aus \(L_\theta\)
 liegt bereits in einer Stufe \(L_{\alpha_n}\).
 Ist eine der erfassten Existenzformeln in \(L\) wahr, so
 hat sie einen Zeugen in \(L_{\alpha_{n+1}}\subseteq L_\theta\).
 Formelinduktion liefert deshalb die behauptete Übereinstimmung.}}
 {\BXLVIIIWhy{Atomare Absolutheit, Junktoren, Zeugen}}
\end{tabproof}

\FormulaThmDeltaK[Die ZF-Axiome gelten in \(L\)]
{\BXLVIIIZFC\vdash\psi^L\quad
 \text{für jedes Axiom }\psi\text{ von }\BXLVIIIZF}
{B48LZF}{
  \DeltaRow{Lesart}{\text{Axiomenschema, keine universelle Wahrheitsdefinition}}
}
\begin{tabproof}
 \proofstep{}{\BXLVIIIText{Extensionalität gilt in der transitiven
 Klasse \(L\): Alle Elemente zweier ihrer Mengen liegen wieder
 in \(L\). Fundierung wird aus dem umgebenden Universum geerbt.
 Die leere Menge und \(\omega\) liegen in \(L\), also gilt auch
 Unendlichkeit.}}
 {\FormulaRefAuto{B48LStages}[thm]}
 \proofstep{}{\BXLVIIIText{Für \(a,b\in L_\alpha\) ist
 \(\{a,b\}\) eine definierbare Teilmenge einer hinreichend großen
 Stufe. Für \(a\in L_\alpha\) ist \(\bigcup a\subseteq L_\alpha\)
 dort durch \(\exists y\in a\ (z\in y)\) definierbar.
 Damit gelten Paarung und Vereinigung.}}
 {\BXLVIIIWhy{Definition der Nachfolgerstufen}}
 \proofstep{}{\BXLVIIIText{Aussonderung für eine feste Formel
 \(\varphi\): Wähle durch Reflexion eine Stufe, welche die Menge
 \(a\), ihre Parameter und die Wahrheit von \(\varphi\) erfasst.
 Die gesuchte Teilmenge
 \(\{x\in a:\varphi^L(x,\bar p)\}\) ist dann über dieser Stufe
 definierbar und liegt in der nächsten.}}
 {\FormulaRefAuto{B48LReflection}[thm]}
 \proofstep{}{\BXLVIIIText{Potenzmenge: Für \(a\in L\) ist
 \(U=\mathcal P(a)\cap L\) im umgebenden Universum eine Menge.
 Ersetzung begrenzt die kleinsten Eintrittsstufen aller
 \(u\in U\). Wähle \(L_\beta\), das \(a\) und alle diese \(u\)
 enthält. Dann
 \(U=\{u\in L_\beta:u\subseteq a\}\) ist über \(L_\beta\)
 definierbar, also selbst in \(L\).}}
 {\BXLVIIIWhy{Äußere Potenzmenge und Ersetzung; beschränkte Absolutheit}}
 \proofstep{}{\BXLVIIIText{Ersetzung für eine feste Formel:
 Sei in \(L\) jedem \(x\in a\) genau ein \(y\) mit
 \(\varphi^L(x,y,\bar p)\) zugeordnet.
 Äußere Ersetzung liefert die Menge aller Werte und eine
 Schranke ihrer Eintrittsstufen. Reflektiere danach
 \(\varphi\) in eine noch größere Stufe \(L_\theta\).
 Die Wertemenge ist dort durch
 \(\exists x\in a\,\varphi(x,y,\bar p)\) definierbar und
 liegt in \(L_{\theta+1}\).}}
 {\BXLVIIIWhy{Eindeutigkeit, äußere Ersetzung und Reflexion}}
 \proofstep{}{\BXLVIIIText{Jeder Schritt ist für die jeweilige
 feste Formel ein ZFC-Beweis. Damit sind alle ZF-Axiome und jede
 einzelne Instanz der beiden Schemata relativiert bewiesen.}}
 {\BXLVIIIWhy{Zusammenfassung der Axiome}}
\end{tabproof}

\subsection{Absolutheit, Auswahl und Kondensation}

Für die Kondensation benötigen wir einen präzisen endlichen
Ausschnitt der bisher bewiesenen Theorie. Mit \(S_L\) bezeichnen wir
die folgende \emph{endliche Liste von Sätzen}, nicht zusätzliche
Grundaxiome: Paarung und Vereinigung; Existenz von \(\omega\);
Existenz der Erfüllungsrelation jeder Mengenstruktur;
Existenz der zu jedem Formelcode und endlichen Parametertupel
gehörigen definierbaren Teilmenge jeder Mengenstruktur;
Existenz der Folge \((L_\xi)_{\xi\le\alpha}\) für jede Ordinalzahl
\(\alpha\). Bei der vierten und fünften Aussage werden Formelcodes
und Parameter als Mengenvariablen quantifiziert; es sind einzelne
Sätze mit der Erfüllungsrelation als definierter Mengenrelation,
keine unendlichen Schemata. Die Hierarchieklausel verlangt auch,
dass jede durch einen Code und Parameter beschriebene Teilmenge
in der jeweiligen Def-Stufe vorkommt.

Alle Sätze dieser Liste sind in ZF bewiesen: durch die
Erfüllungskonstruktion, Aussonderung für die eine Erfüllungsformel,
Ersetzung über die Codes und Parameter sowie transfinite Rekursion.
\(\omega\) kann dabei als kleinste von \(0\) verschiedene
Limesordinalzahl charakterisiert werden.

\FormulaThmDeltaK[Absolutheit der konstruktiblen Hierarchie]
{\begin{gathered}
 N\text{ transitiv},\quad N\models S_L,\quad \alpha\in N\cap\BXLVIIIOrd\\
 \vdash\ (L_\alpha)^N=L_\alpha
 \end{gathered}}
{B48LAbsolute}{
  \DeltaRow{Struktur}{(N,\in)}
}
\begin{tabproof}
 \proofstep{1}{N\text{ transitiv},\quad N\models S_L}{\rA}
 \proofstep{1}{\BXLVIIIText{Ordinalsein ist in einer transitiven
 Menge absolut: Die beschränkten Bedingungen Transitivität und
 strikte \(\in\)-Linearität sind absolut, Fundierung gilt außen.
 Die intern kleinste nichtnullige Limesordinalzahl ist deshalb
 das wirkliche \(\omega\).}}
 {\BXLVIIIWhy{Transitivität und Fundierung}}
 \proofstep{1}{\BXLVIIIText{Durch Paarung und Vereinigung enthält
 \(N\) alle wirklichen endlichen Tupel seiner Elemente.
 Die Formelcodes und ihre endlichen Aufbauten sind daher dieselben
 wie außen. Für eine Struktur \(X\in N\) quantifizieren ihre
 Erfüllungsklauseln über genau dieselben Elemente von \(X\).
 Formelinduktion zeigt die Absolutheit dieser Erfüllung.}}
 {\BXLVIIIWhy{\(\omega^N=\omega\); Erfüllungsrelation aus \(S_L\)}}
 \proofstep{1}{\BXLVIIIText{Die Def-Existenzklausel aus \(S_L\)
 liefert zu jedem wirklichen Code und jedem Parametertupel aus
 \(X\) die entsprechende Teilmenge in \(N\).
 Da ihre Erfüllung absolut ist, ist
 \(\BXLVIIIDef(X)^N=\BXLVIIIDef(X)\).
 Es fehlen also auch keine außerhalb beschreibbaren Teilmengen.}}
 {\BXLVIIIWhy{Vollständige Erfassung der Codes und Parameter}}
 \proofstep{1}{\BXLVIIIText{Die in \(N\) vorhandene Hierarchiefolge
 bis \(\alpha\) stimmt nun durch transfinite Induktion mit der
 äußeren überein: leere Anfangsstufe, absoluter Def-Schritt,
 an Limiten dieselbe Vereinigung.}}
 {\BXLVIIIWhy{Hierarchieexistenz und Eindeutigkeit der Rekursion}}
\end{tabproof}

\noindent
Der Beweis gilt ebenso schematisch für die transitive Klasse \(L\).
Da sie ZF erfüllt, konstruiert sie intern genau dieselben \(L\)-Stufen.
Jedes ihrer Elemente liegt in einer solchen Stufe. Folglich
\[
 (V=L)^L.
\]

\FormulaThmDeltaK[Auswahl im konstruktiblen Universum]
{\BXLVIIIZFC\vdash\mathsf{AC}^L}
{B48LChoice}{
  \DeltaRow{Bereits bewiesen}{L\models\BXLVIIIZF,\quad (V=L)^L}
}
\begin{tabproof}
 \proofstep{}{\BXLVIIIText{Arbeite intern in ZF mit \(V=L\).
 Konstruiere gleichzeitig mit den Stufen Wohlordnungen:
 Die alten Elemente kommen zuerst. Neue Elemente einer
 Nachfolgerstufe werden durch ihre kleinsten Definitionscodes
 geordnet.}}
 {\BXLVIIIWhy{Transfinite Rekursion}}
 \proofstep{}{\BXLVIIIText{Ein Definitionscode besteht aus einer
 natürlichen Formelnummer und einem endlichen Parametertupel.
 Ordne erst die Formelnummern, dann die Tupel fester Länge
 lexikographisch nach der bereits vorhandenen Wohlordnung.
 Jede definierbare Teilmenge erhält ihren kleinsten Code.
 Verschiedene Mengen haben verschiedene kleinste Codes.}}
 {\BXLVIIIWhy{Wohlordnung endlicher lexikographischer Produkte}}
 \proofstep{}{\BXLVIIIText{Die neue Wohlordnung erweitert die
 alte als Anfangsstück. An Limiten wird vereinigt.
 Eine nichtleere Teilmenge der Vereinigung besitzt ein Element
 in einer frühesten Eintrittsstufe und dort ein kleinstes;
 die Vereinigung ist also wieder eine Wohlordnung.}}
 {\BXLVIIIWhy{Ordinalwohlordnung}}
 \proofstep{}{\BXLVIIIText{Wegen \(V=L\) liegt jede Menge in einer
 Stufe. Die Einschränkung ihrer Stufenwohlordnung wohlordnet
 diese Menge. Für jede Familie nichtleerer Mengen wählt man
 damit jeweils das kleinste Element ihrer Vereinigung.}}
 {\BXLVIIIWhy{Aussonderung und Ersetzung; Wohlordnung impliziert Auswahl}}
 \proofstep{}{\mathsf{AC}^L}
 {\BXLVIIIWhy{ZF und \(V=L\) gelten in \(L\)}}
\end{tabproof}

\FormulaThmDeltaK[Kondensation für den benötigten endlichen Ausschnitt]
{\begin{gathered}
 N\text{ transitiv},\quad N\models S_L+(V=L)\\
 \vdash\ N=L_\delta,\qquad \delta=N\cap\BXLVIIIOrd
 \end{gathered}}
{B48Condensation}{
  \DeltaRow{Menge}{N}
}
\begin{tabproof}
 \proofstep{1}{N\text{ transitiv},\quad N\models S_L+(V=L)}{\rA}
 \proofstep{1}{\BXLVIIIText{Der Ordinalteil \(\delta\) ist eine
 Ordinalzahl. Er ist eine Limesordinalzahl, denn zu jedem
 \(\alpha\in N\) liegt durch Paarung und Vereinigung auch
 \(\alpha+1\) in \(N\); außerdem enthält \(N\) die Menge \(\omega\).}}
 {\BXLVIIIWhy{Transitivität und \(S_L\)}}
 \proofstep{1}{\BXLVIIIText{Für jedes \(\alpha<\delta\) ist die
 interne Menge \((L_\alpha)^N\) vorhanden und gleich \(L_\alpha\).
 Transitivität liefert \(L_\alpha\subseteq N\).
 Also \(L_\delta\subseteq N\).}}
 {\FormulaRefAuto{B48LAbsolute}[thm]}
 \proofstep{1}{\BXLVIIIText{Für jedes \(x\in N\) liefert das
 interne \(V=L\) eine Ordinalzahl \(\alpha\in N\) mit
 \(x\in(L_\alpha)^N=L_\alpha\).
 Daher \(N\subseteq L_\delta\).}}
 {\BXLVIIIWhy{Absolutheit und internes \(V=L\)}}
 \proofstep{1}{N=L_\delta}{\BXLVIIIWhy{Beide Inklusionen}}
\end{tabproof}

\FormulaThmDeltaK[Frühes Auftreten von Teilmengen]
{\BXLVIIIZFC+(V=L)\vdash\
 \forall\kappa\ge\aleph_0\ \forall x\subseteq\kappa\
 (x\in L_{\kappa^+})}
{B48EarlySubsets}{
  \DeltaRow{Kardinalzahl}{\kappa}
}
\begin{tabproof}
 \proofstep{1}{V=L,\quad\kappa\text{ unendlich},\quad x\subseteq\kappa}
 {\rA}
 \proofstep{1}{\BXLVIIIText{Wähle durch Reflexion eine hinreichend
 große Stufe \(L_\theta\), die \(x,\kappa\) enthält und
 \(S_L+(V=L)\) erfüllt. Die Liste ist endlich und alle ihre
 Sätze gelten unter \(V=L\).}}
 {\FormulaRefAuto{B48LReflection}[thm]}
 \proofstep{1}{\kappa\cup\{\kappa,x\}\subseteq X\prec L_\theta,
 \qquad |X|=\kappa}
 {\FormulaRefAuto{B48SkolemHull}[thm]}
 \proofstep{1}{\BXLVIIIText{Die Relation \(\in\) auf \(X\) ist
 fundiert. Sie ist extensional, weil die Elementarität zu
 verschiedenen Mengen aus \(X\) einen unterscheidenden
 Mitgliedschaftszeugen in \(X\) liefert.
 Ihr Kollaps ist eine transitive Menge \(N\), und
 \(N\models S_L+(V=L)\).}}
 {\FormulaRefAuto{B48Mostowski}[thm]}
 \proofstep{1}{\BXLVIIIText{Der Kollaps \(\pi\) fixiert jede
 Ordinalzahl unterhalb \(\kappa\), nach Induktion, weil alle
 deren Elemente in \(X\) liegen. Da auch \(\kappa\in X\),
 gilt \(\pi(\kappa)=\kappa\). Weil \(x\subseteq\kappa\) und
 \(x\in X\), folgt ebenso \(\pi(x)=x\).}}
 {\BXLVIIIWhy{Kollapsgleichung; vollständige Aufnahme von \(\kappa\)}}
 \proofstep{1}{N=L_\delta,\quad x\in L_\delta,\quad
 |\delta|\le|N|=\kappa}
 {\FormulaRefAuto{B48Condensation}[thm]}
 \proofstep{1}{\delta<\kappa^+,\qquad x\in L_{\kappa^+}}
 {\BXLVIIIWhy{Definition von \(\kappa^+\); Monotonie der Hierarchie}}
\end{tabproof}

\subsection{GCH und relative Konsistenz}

\FormulaDefDeltaK[Verallgemeinerte Kontinuumshypothese]
{\BXLVIIIGCH:\Longleftrightarrow\
 \forall\kappa\text{ unendliche Kardinalzahl}\ (2^\kappa=\kappa^+)}
{B48GCHDef}{
  \DeltaRow{Nachfolgerkardinal}{\kappa^+}
}

\FormulaThmDeltaK[Gödels relative Konsistenz von GCH]
{\BXLVIIICon(\BXLVIIIZFC)\ \Longrightarrow\
 \BXLVIIICon(\BXLVIIIZFC+\BXLVIIIGCH)}
{B48GodelConsistency}{
  \DeltaRow{Metatheoretische Voraussetzung}{\BXLVIIICon(\BXLVIIIZFC)}
}
\begin{tabproof}
 \proofstep{}{\BXLVIIIText{In ZFC mit \(V=L\) gilt
 \(\mathcal P(\kappa)\subseteq L_{\kappa^+}\).
 Für unendliches \(\kappa\) folgt
 \(2^\kappa\le|L_{\kappa^+}|=\kappa^+\).}}
 {\BXLVIIIWhy{Frühes Auftreten und Größe der \(L\)-Stufen}}
 \proofstep{}{\BXLVIIIText{Die Diagonalmenge
 \(D=\{\xi<\kappa:\xi\notin f(\xi)\}\) widerspricht jeder
 behaupteten Surjektion \(f:\kappa\to\mathcal P(\kappa)\).
 Zusammen mit der Singleton-Injektion folgt
 \(\kappa<2^\kappa\), also \(\kappa^+\le2^\kappa\).}}
 {\BXLVIIIWhy{Cantors Diagonalargument; ausführlich in Kapitel 4}}
 \proofstep{}{\BXLVIIIText{Somit beweist ZFC mit \(V=L\) die GCH.
 Da \(L\) ZFC und intern \(V=L\) erfüllt, ist in ZFC jede
 ZFC-Axiomrelativierung auf \(L\) und außerdem
 \(\BXLVIIIGCH^L\) bewiesen.}}
 {\BXLVIIIWhy{ZF, Auswahl und Absolutheit für \(L\)}}
 \proofstep{4}{\BXLVIIICon(\BXLVIIIZFC)}{\rA}
 \proofstep{4}{M\models\BXLVIIIZFC,\quad |M|\le\aleph_0}
 {\FormulaRefAuto{B48Completeness}[thm]}
 \proofstep{4}{\BXLVIIIText{Beschränke die Grundmenge von \(M\)
 auf die Elemente, für die \(M\) die Formel \(x\in L\) erfüllt,
 und die Relation entsprechend. Nenne diese äußere
 Mengenstruktur \(L^M\). Formelinduktion ergibt
 \(L^M\models\psi\) genau dann, wenn \(M\models\psi^L\).}}
 {\BXLVIIIWhy{Relativierung; keine äußere Fundiertheit von \(M\) nötig}}
 \proofstep{4}{L^M\models\BXLVIIIZFC+\BXLVIIIGCH}
 {\BXLVIIIWhy{Korrektheit der in Zeile 3 bewiesenen einzelnen Sätze}}
 \proofstep{4}{\BXLVIIICon(\BXLVIIIZFC+\BXLVIIIGCH)}
 {\FormulaRefAuto{B48ModelConsistency}[thm]}
\end{tabproof}

\noindent
Die Konstruktion benutzt nicht die Behauptung, \(L^M\) sei von außen
transitiv. Entscheidend ist die Interpretation jeder festen Formel.
Damit reicht tatsächlich die Konsistenzannahme, nicht die stärkere
Annahme eines transitiven Modells.
